\documentclass{article}

\title{Ch. 14 HW \#1 Solutions}
\author{Principles of Microeconomics}
\date{\today}

\begin{document}

\maketitle

\subsection*{Quick Quiz}

\begin{enumerate}

\item (c) -- His explicit cost is \$100 for seeds. His implicit cost is \$200 from the income he forewent by not teaching banjo lessons.

\item (a) -- His explicit cost is \$10 for ingredients. His implicit cost is \$40 from the income he forewent by not mowing his neighbor's lawn. That means his accounting profit is $60 - 10 = 50$, and his economic profit is $60 - (10 + 40) = 10$. 

\end{enumerate}

\subsection*{Questions for Review}

\begin{enumerate}

\item Profit = Total Revenue - Total Cost

\item The money that firms spend establishing their business could earn interest if the firm saved it instead. We call the interest that they would have earned the cost of capital. Since the cost of capital in this case is an implicit cost -- it doesn't require a cash outlay -- accountants would ignore it.
\end{enumerate}

\end{document}